% Reporte de avances de tesis para el comite tutor (PCIC, UNAM), agosto-noviembre de 2026.
% Compilar: latexmk reporte_avances.tex  (latexmkrc elige LuaLaTeX y agrega ../presentacion).
% Las cifras vienen de docs/progreso.md, docs/bitacora_decisiones.md, docs/resultados_*.md,
% docs/traspaso_tinto.md y los CSV citados en los comentarios de cada figura
% (presentacion/figuras/). Las graficas se reutilizan tal cual de la presentacion para que
% los numeros sean los mismos en los dos documentos.
% Version 2 (2026-11-24): ampliada a pedido del tesista.
\documentclass[11pt, letterpaper]{article}

\usepackage[spanish, es-nodecimaldot, es-noshorthands, es-tabla]{babel}
\usepackage{fontspec}
\usepackage[margin=2.5cm]{geometry}
\usepackage{csquotes}
\usepackage[backend=biber, style=ieee]{biblatex}
\addbibresource{../docs/referencias.bib}
\usepackage{amsmath}
\usepackage{booktabs}
\usepackage{tabularx}
\usepackage{colortbl}
\usepackage{enumitem}
\usepackage{adjustbox}   % \adjustbox{max width=\linewidth}: encoge una figura solo si no cabe
\usepackage{colores-tesis}  % paleta y estilos de las graficas (de ../presentacion)
\usepackage{iconos-tesis}   % iconos de la figura de contexto (de ../presentacion)
\usepackage{tikz}
\usetikzlibrary{arrows.meta, calc, positioning}
\usepgfplotslibrary{groupplots}
\usepackage{caption}
\captionsetup{font=small, labelfont=bf}
\usepackage{float}
\usepackage[hidelinks]{hyperref}

\newcolumntype{Y}{>{\raggedright\arraybackslash}X}
\newcolumntype{C}{>{\centering\arraybackslash}X}
\setlist{itemsep=2pt, topsep=4pt}
% Dato que falta llenar a mano (numero de cuenta, semestre, comite, fechas)
\newcommand{\porllenar}[1]{\textcolor{acento}{[#1]}}
% Figura de la presentacion, encogida al ancho del texto si hace falta
\newcommand{\figurapresentacion}[1]{%
  \adjustbox{max width=\linewidth}{\input{../presentacion/figuras/#1}}}
% Palomita y tache (los mismos de la presentacion)
\newcommand{\marcasi}{\tikz[baseline=-0.6ex]\draw[ensX, line width=1.6pt,
  line cap=round, line join=round] (-0.12,0) -- (-0.03,-0.09) -- (0.13,0.1);}
\newcommand{\marcano}{\tikz[baseline=-0.6ex]\draw[tenue, line width=1.3pt,
  line cap=round] (-0.09,-0.09) -- (0.09,0.09) (-0.09,0.09) -- (0.09,-0.09);}

\begin{document}

% =====================================================================
\begin{center}
  {\large\bfseries Posgrado en Ciencia e Ingeniería de la Computación · UNAM}\\[4pt]
  {\Large\bfseries Reporte de avances de tesis de maestría}\\[2pt]
  {\large Periodo: agosto a noviembre de 2026}
\end{center}

\medskip
\noindent
\begin{tabularx}{\linewidth}{@{}l Y@{}}
  \toprule
  \textbf{Alumno} & Ing. José Emiliano Lili Beltrán \\
  \textbf{Número de cuenta} & \porllenar{por llenar} \\
  \textbf{Semestre} & \porllenar{por llenar} \\
  \textbf{Tutor principal} & Dr. José Jaime Camacho Escoto \\
  \textbf{Comité tutor} & \porllenar{por llenar} \\
  \textbf{Título de trabajo} & Detección de anomalías en redes de monitoreo atmosférico:
    aprendizaje profundo con contraste espacio-temporal y exploración con imágenes
    sintéticas TINTO \\
  \textbf{Fecha} & 24 de noviembre de 2026 \\
  \bottomrule
\end{tabularx}

\tableofcontents

% =====================================================================
\section{Resumen del periodo}

Durante el periodo se completaron las dos etapas experimentales de la tesis y se cerró la
elección del detector final. El trabajo partió de un ataque concreto a la capa de
comunicación (\emph{bit flipping} en LoRaWAN~1.1) y, por acuerdo con el tutor, se amplió a
un problema más general: detectar, desde el servidor central, qué estación de una red de
monitoreo reporta datos anómalos, sin importar si la causa es una falla, una
descalibración o un ataque. Los avances principales son:

\begin{enumerate}
  \item \textbf{Etapa 1, cerrada:} caracterización del ataque de \emph{bit flipping} en
    LoRaWAN~1.1 sobre mediciones reales de ozono de la Red Automática de Monitoreo
    Atmosférico de la Ciudad de México (RAMA): qué posiciones de bit producen un daño real
    (cambio de la banda de calidad del aire) y cuáles detecta un modelo LSTM.
  \item \textbf{Reorientación acordada con el tutor} (octubre de 2026): de un ataque a un
    solo sensor a la \textbf{detección de estaciones anómalas a nivel de red}.
  \item \textbf{Etapa 2:} diseño de un protocolo experimental con separación temporal
    estricta (2020--2026) y desarrollo de un detector propio, \textbf{NERAD}
    (\emph{Neighbour-filtered Ensemble Recurrent Anomaly Detection}): un ensamble de
    5~redes LSTM que compara cada estación con su propia historia y con sus 4~vecinas más
    cercanas, y que descarta a las vecinas sospechosas.
  \item \textbf{Comparación con el estado del arte:} NERAD supera a PARAAD-SA
    \cite{allka2024paraad} en el banco de pruebas de la RAMA con menos de la mitad de
    falsas alarmas; el método base también lo supera en datos de Barcelona (validación
    externa).
  \item \textbf{Pruebas de estrés} ante ataques coordinados a las vecinas y sesgos lentos,
    que delimitan con números dónde funciona el detector y dónde no.
  \item \textbf{Exploración con imágenes sintéticas TINTO}: pilotos cerrados y una primera
    comparación pareada con el detector final.
  \item \textbf{Material de difusión:} presentación de avances en \LaTeX{} (Beamer) con
    todas las gráficas generadas a partir de los resultados guardados.
\end{enumerate}

% =====================================================================
\section{Planteamiento}

\subsection{Contexto y problema}

La RAMA mide ozono (O$_{\text{3}}$) cada hora en decenas de estaciones de la Zona
Metropolitana del Valle de México. Con esos datos se informa a la población la calidad del
aire y se activan, cuando corresponde, las fases de contingencia ambiental. La
concentración de ozono varía mucho a lo largo del día (es baja en la madrugada y alcanza
su máximo por la tarde) y entre zonas de la ciudad.

Cuando un sensor falla, se descalibra o su mensaje se altera durante la transmisión, el
dato que llega al servidor central suele tener un formato válido y un valor
\emph{físicamente posible}. Por eso un umbral fijo o una verificación de rango no lo
detectan. La pregunta práctica es cómo dar al operador de la red un aviso, al cierre de
cada día, de qué estaciones requieren revisión, sin instalar sensores redundantes y sin
saturarlo de falsas alarmas (Figura~\ref{fig:contexto}).

\begin{figure}[H]
  \centering
  \figurapresentacion{contexto_problema}
  \caption{Contexto del problema: (1) una red de estaciones que reportan ozono cada hora;
    (2) un sensor o un mensaje alterado que entrega un valor falso con formato válido;
    (3) el reto operativo de avisar al cierre del día qué estación debe revisarse.}
  \label{fig:contexto}
\end{figure}

\subsection{Hipótesis}

Es posible detectar, con un alto grado de confiabilidad, anomalías en la información
reportada en redes de sensores utilizando el historial reciente de cada estación y el de
sus estaciones vecinas, sin incurrir en un exceso de falsos positivos.

\subsection{Objetivos}

\paragraph{Objetivo general.} Diseñar un detector de anomalías basado en modelos de
aprendizaje profundo que combine la historia temporal de cada estación con la de sus
estaciones vecinas, con alta capacidad de detección y una tasa de falsas alarmas acotada,
en una red de sensores de ozono.

\paragraph{Objetivos específicos.}
\begin{enumerate}[label=OE\arabic*.]
  \item Caracterizar el daño y la detectabilidad del \emph{bit flipping} en LoRaWAN~1.1
    según la posición de bit atacada.
  \item Construir un banco de pruebas reproducible con anomalías sintéticas inyectadas
    sobre datos reales de la RAMA y una separación temporal entrenamiento / selección /
    calibración / prueba.
  \item Desarrollar un detector supervisado que use la historia propia y la de las
    vecinas, y compararlo con detectores del mismo tipo y con un método publicado.
  \item Delimitar el alcance del detector ante ataques coordinados y sesgos de largo plazo.
  \item Explorar una representación alternativa de los datos como imágenes (TINTO).
\end{enumerate}

\subsection{Avance por objetivo}

\begin{center}
\small
\renewcommand{\arraystretch}{1.25}
\begin{tabularx}{\linewidth}{@{}l l Y@{}}
  \toprule
  \textbf{Objetivo} & \textbf{Estado} & \textbf{Evidencia principal} \\
  \midrule
  OE1 & Cumplido & Detección por bit del LSTM V4 en 2024--2026; zona explotable en los
    bits 3 a 5 (sección~\ref{sec:etapa1}). \\
  OE2 & Cumplido & Protocolo con más de 25~revisiones fechadas; banco 2024 a julio de 2026
    con 30~calendarios de ataque (sección~\ref{sec:protocolo}). \\
  OE3 & Cumplido & NERAD contra PARAAD-SA, V4 supervisado y regla de vecinas
    (secciones \ref{sec:evolucion} y~\ref{sec:resultados}). \\
  OE4 & Cumplido & Ataque coordinado, sesgo sostenido y fusión
    (sección~\ref{sec:estres}). \\
  OE5 & En curso & Pilotos T1 y T2 cerrados; comparación pareada inicial
    (sección~\ref{sec:tinto}). \\
  \bottomrule
\end{tabularx}
\end{center}

% =====================================================================
\section{Trabajos relacionados}

Se revisó la literatura sobre detección de anomalías en redes de sensores ambientales y
sobre detección de intrusiones en IoT. Los tres trabajos más cercanos se resumen en la
Tabla~\ref{tab:relacionados}.

\begin{table}[H]
\centering
\small
\caption{Trabajos relacionados y diferencias con este trabajo. Las cifras son las
  reportadas por cada artículo y no son comparables entre sí.}
\label{tab:relacionados}
\renewcommand{\arraystretch}{1.6}
\setlength{\tabcolsep}{5pt}
\begin{tabularx}{\linewidth}{>{\bfseries}p{2.3cm} C C C >{\columncolor{panel}}C}
  \toprule
  & \textbf{PARAAD-SA}~\cite{allka2024paraad}\newline{\footnotesize\color{tenue}Autoencoder recurrente}
  & \textbf{STC-RAAD}~\cite{allka2024stcraad}\newline{\footnotesize\color{tenue}Autoencoder BiLSTM}
  & \textbf{FedWLA}~\cite{garciasaez2025fedwla}\newline{\footnotesize\color{tenue}Árboles federados}
  & \textcolor{ensX}{\textbf{NERAD}}\newline{\footnotesize\color{tenue}Este trabajo: 5 LSTM} \\
  \midrule
  Datos        & O$_{\text{3}}$, sensores & O$_{\text{3}}$, sensores & Tráfico de red & O$_{\text{3}}$, sensores \\
  Supervisado  & \marcano & \marcano & \marcasi & \marcasi \\
  Usa vecinas  & \marcano & \marcasi & {\color{tenue}---} & \marcasi\ {\footnotesize con filtro} \\
  Resultado    & R=0.81 · P=0.84 & R=0.82--0.96 · P=0.85--0.91 & F1=0.98--0.99 & Sección~\ref{sec:resultados} \\
  \bottomrule
\end{tabularx}
\end{table}

\begin{itemize}
  \item \textbf{PARAAD-SA} \cite{allka2024paraad} (UPC Barcelona, IEEE TNSE 2024) detecta
    anomalías de sesgo y deriva en ozono de 8~estaciones de Barcelona con un autoencoder
    recurrente sin supervisión. Es nuestro comparador principal, por pedido del tutor y
    porque es el método de red del mismo grupo.
  \item \textbf{STC-RAAD} \cite{allka2024stcraad} (mismo grupo, IEEE IoT Journal 2024) usa a
    los sensores vecinos mediante un grafo. Es el trabajo conceptualmente más cercano al
    nuestro. Sin embargo, se ajustó con 25\,\% de días anómalos, vecinas siempre limpias y
    un umbral elegido con etiquetas, y no tiene código público, por lo que su comparación
    directa queda pendiente.
  \item \textbf{FedWLA} \cite{garciasaez2025fedwla} (UCLM y U. de Zaragoza, IEEE SMC 2025)
    es un sistema federado de detección de intrusiones. Ataca la red, no el dato, y evalúa
    con 28--49\,\% de ataques; en nuestro banco sólo 0.77\,\% de los estación-días están
    atacados, lo que explica sus métricas tan altas.
\end{itemize}

\textbf{Diferencia de este trabajo:} es supervisado (aprende con anomalías sintéticas
etiquetadas), usa a las vecinas \emph{y} filtra a las sospechosas, y se evalúa en cuatro
escenarios de ataque sobre el dato con una proporción realista de días atacados.

% =====================================================================
\section{Actividades realizadas}

\subsection{Datos y preparación (agosto de 2026)}

\begin{itemize}
  \item Inventario y auditoría de las series horarias de ozono de la RAMA de 2020 a
    julio de 2026. El formato es estable en todos los años; cada archivo tiene su suma de
    verificación. Se evalúan 31~estaciones: HGM sólo sirve de contexto, XAL está
    inactiva y COY, SJA y SFE quedan fuera por falta de datos.
  \item Datos complementarios: series de 2000 a 2019 (sólo como fondo de entrenamiento),
    datos validados del SIMAT para agosto--septiembre de 2026 y ozono de 13~estaciones de
    Barcelona (2014--2019) para la validación externa.
  \item Función de decisión normativa: bandas de calidad del aire de la
    NOM-172-SEMARNAT-2023 (Tabla~\ref{tab:nom}). Umbral experimental de contingencia:
    155~ppb, aplicado igual en todos los años.
  \item Simulación del transporte LoRaWAN~1.1: codificación CayenneLPP y cifrado AES en
    modo contador, sobre el que se aplican las alteraciones.
\end{itemize}

\begin{table}[H]
  \centering
  \small
  \caption{Bandas de calidad del aire para ozono (promedio horario), NOM-172-SEMARNAT-2023.}
  \label{tab:nom}
  \begin{tabular}{@{}lc@{}}
    \toprule
    \textbf{Banda} & \textbf{Concentración (ppb)} \\
    \midrule
    Buena & $<58$ \\
    Aceptable & 58--90 \\
    Mala & 90--135 \\
    Muy mala & 135--175 \\
    Extremadamente mala & $>175$ \\
    \bottomrule
  \end{tabular}
\end{table}

\subsection{Etapa 1: \emph{bit flipping} en LoRaWAN 1.1 (agosto a octubre de 2026)}
\label{sec:etapa1}

\paragraph{El ataque.} En LoRaWAN la confidencialidad y la integridad se protegen con
mecanismos independientes. La carga útil se cifra con AES en modo contador, que es un XOR
contra un flujo de clave: invertir un bit del texto cifrado invierte el mismo bit del texto
claro, sin necesidad de conocer la clave. La integridad la protege el MIC, pero su
verificación termina en el servidor de red; el tramo entre el servidor de red y el de
aplicación queda sin protección de extremo a extremo. Es una vulnerabilidad de
arquitectura que afecta igual a las versiones 1.0.x y~1.1.

\paragraph{Medida del daño.} El daño no se midió como el tamaño del desplazamiento, sino
como un \emph{cambio de decisión}: hay daño si la banda normativa calculada con el valor
recibido es distinta de la calculada con el valor real,
\[
  D = \mathbf{1}\bigl[f(x_{\text{real}}) \neq f(x_{\text{recibido}})\bigr].
\]
Un salto de 32~ppb lejos de un corte no produce daño; uno de 2~ppb junto al corte cambia
la banda que se comunica a la población.

\paragraph{Resultados.}
\begin{itemize}
  \item \textbf{Existe un óptimo interior del atacante.} Las alteraciones pequeñas son
    indetectables pero casi nunca cruzan un umbral; las grandes producen valores fuera de
    lo observado y se detectan fácilmente. El máximo daño no detectado está en medio.
  \item \textbf{El contexto espacial cierra la ventana} (Tabla~\ref{tab:capa1}). Con un
    LSTM que estima el valor esperado de la estación a partir del resto de la red, la
    detección mejora y el óptimo del atacante se desplaza a 16~ppb, donde sólo 11.4\,\%
    de los ataques llega a cambiar la banda.
  \item \textbf{Deriva estacional:} un perfil horario fijo de enero--octubre subestima en
    10.6~ppb la concentración de noviembre--diciembre a las 16~h. Se adoptó un perfil
    causal de 14~días, que empeora un poco el error medio (9.35 contra 8.39~ppb) pero
    reduce el sesgo de +18.7 a $-$1.7~ppb. Este perfil de 14~días se conservó después en
    la etapa~2.
  \item \textbf{Detector multianual (V4):} 31~modelos LSTM, uno por estación, entrenados
    con 2020--2022, calibrados en 2023 y probados en 2024, 2025 y 2026. Reducen el error
    frente al perfil en 30 de 31~estaciones (mediana de reducción de 15.9\,\%). La
    detección depende del bit atacado (Tabla~\ref{tab:bits}).
  \item \textbf{Daño no detectado:} se concentra en los bits 3 a~5. En 2024, el bit~5
    permitió cambiar sin alerta la banda del máximo diario de la red en 85~días (23\,\%).
  \item \textbf{Contramedida probada sin éxito:} una codificación alternativa del campo de
    valor que reparte el peso de los bits 4 y~5. En 27~estaciones, con 100~campañas por
    estación, aumentó los cambios de banda no detectados de 255 a 567 en 179\,399~intentos.
    Se documenta como resultado negativo.
\end{itemize}

\begin{table}[H]
  \centering
  \small
  \caption{Detección del \emph{bit flipping} sin y con contexto espacial (estación CCA,
    2025).}
  \label{tab:capa1}
  \begin{tabular}{@{}lcc@{}}
    \toprule
    & \textbf{Medición aislada} & \textbf{Con contexto espacial (LSTM)} \\
    \midrule
    Detección, desplazamiento de 32~ppb & 41.4\,\% & 84.6\,\% \\
    Detección, desplazamiento de 64~ppb & 72.9\,\% & 98.7\,\% \\
    Máximo daño no detectado & 27.1\,\% (64~ppb) & 9.1\,\% (16~ppb) \\
    \bottomrule
  \end{tabular}
\end{table}

\begin{table}[H]
  \centering
  \small
  \caption{Detección aproximada del V4 multianual según el bit atacado (grupo principal,
    probabilidad de ataque de 5\,\% por lectura, 2024--2026). Las falsas alarmas en prueba
    fueron de 6.6 a 7.3\,\%, frente al 5\,\% calibrado.}
  \label{tab:bits}
  \begin{tabular}{@{}lccccc@{}}
    \toprule
    \textbf{Bit} & 0--2 & 4 & 5 & 6 & 7 \\
    \midrule
    \textbf{Recall} & $\approx$7\,\% & 20--22\,\% & 79--80\,\% & $>$99\,\% & 100\,\% \\
    \bottomrule
  \end{tabular}
\end{table}

Con el tutor se concluyó que el detector estaba topado para este ataque: más años de
datos, vecinas de la misma hora o vecinas cercanas apenas movieron el resultado. La etapa
queda como primera contribución de la tesis: una caracterización formal de qué bits
producen daño y cuáles detecta un LSTM.

\subsection{Reorientación (octubre de 2026)}

En la junta del 5 de octubre se evaluó un segundo ataque a LoRaWAN, el drenado de energía
por interferencia en las ventanas de recepción. En la junta del 12 de octubre se descartó
por el costo de montar un banco físico, y se acordó la línea vigente, sugerida por el
Dr.~Orozco (UCLM): dejar de ver el ataque a \emph{un} sensor y ver la \emph{red como un
todo}, suponiendo que un nodo ya está descompuesto o es malicioso y usando a los demás
para detectarlo. Los acuerdos fueron:

\begin{itemize}
  \item La anomalía no tiene por qué ser maliciosa: puede ser falla o descalibración.
  \item Datos: matriz con una fila por hora y una columna por estación de la RAMA.
  \item Anomalías sintéticas de magnitud variable, de 1 a 7~días, en estaciones al azar.
  \item División por años completos, del modelo más simple al más complejo, con la meta
    de una detección del orden de 80--90\,\%.
\end{itemize}

\subsection{Etapa 2: protocolo experimental (octubre de 2026)}
\label{sec:protocolo}

Antes de programar se escribió un protocolo con las reglas que no cambian. Cada cambio
posterior se registró con fecha \emph{antes} de ver resultados; a la fecha hay más de
25~revisiones documentadas.

\paragraph{Anomalía.} En un episodio, el valor $x$ de la estación atacada se sustituye por
\[
  x' = x\,(1+\delta) \quad\text{(aumento)} \qquad\text{o}\qquad
  x' = \frac{x}{1+\delta} \quad\text{(reducción)},
\]
con el sentido sorteado por episodio. Los episodios duran de 1 a 7~días completos y
empiezan a las 00~h. Sólo se alteran las 31~estaciones objetivo y sólo en horas con dato.
Más adelante se agregaron otros escenarios (Tabla~\ref{tab:escenarios}).

\begin{table}[H]
\centering
\small
\caption{Escenarios de anomalía evaluados (modelo de amenaza).}
\label{tab:escenarios}
\renewcommand{\arraystretch}{1.35}
\begin{tabularx}{\linewidth}{@{}>{\raggedright\bfseries}p{4.3cm} Y >{\raggedright\arraybackslash}p{3cm}@{}}
  \toprule
  Escenario & Magnitudes evaluadas & Duración \\
  \midrule
  1. Anomalía proporcional & $\pm\delta$, con $\delta$ = 20, 50 y 100\,\% & 1 a 7 días \\
  2. Escalón aditivo fijo  & $\pm$10, $\pm$20 y $\pm$40~ppb & 1 a 7 días \\
  3. Ataque coordinado a vecinas & 0, 1, 2 y 4 vecinas alteradas, con las magnitudes de 1 y 2 & 1 a 7 días \\
  4. Sesgo lento o sostenido & +20 y +40~ppb: escalón, rampa de 28~días y deriva de 56~días & 8 semanas \\
  \bottomrule
\end{tabularx}
\end{table}

\paragraph{Calendarios pareados.} Un \emph{calendario} es un sorteo de qué estaciones se
atacan, qué días y en qué sentido. El mismo calendario se aplica a todas las magnitudes,
de modo que las comparaciones entre magnitudes y entre detectores usan exactamente los
mismos días. Las semillas se derivan de 42 de forma reproducible. La etiqueta nunca entra
al modelo como dato.

\paragraph{División temporal.} Cada etapa usa sólo datos anteriores a ella
(Figura~\ref{fig:protocolo}). Nada se ajusta mirando los datos de prueba, y los exámenes
finales se abren una sola vez.

\begin{figure}[H]
  \centering
  \figurapresentacion{protocolo}
  \caption{Protocolo temporal del detector final. Los modelos se entrenan con 2020--2021,
    el diseño se elige en 2022 (y después se reentrena con 2020--2022), el umbral se fija
    en 2023 sin ataques, el banco de pruebas cubre de 2024 a julio de 2026 (30
    calendarios de ataque) y el despliegue, agosto y septiembre de 2026.}
  \label{fig:protocolo}
\end{figure}

\paragraph{Unidad de evaluación.} Al principio se evaluó la \emph{hora de red} (¿hay
alguna anomalía en la red en esta hora?). Desde la revisión~11 se usa la
\emph{estación-día}: ¿esta estación tuvo hoy una anomalía? Es la pregunta que responde el
operador al cierre del día y permite localizar la estación.

\paragraph{Umbrales y métricas.} El umbral de alarma se fija como un percentil del puntaje
en 2023 sin ataques: con p99, se espera que 1 de cada 100 días normales dé alarma. Las
métricas se definen así:
\begin{itemize}
  \item \textbf{Recall} (sensibilidad): de los días atacados, qué fracción dio alarma.
  \item \textbf{Precisión}: de las alarmas, qué fracción correspondía a un día atacado.
  \item \textbf{F1}: media armónica de recall y precisión; resume las dos en un número.
  \item \textbf{Falsas alarmas}: días sin ataque con alarma, expresados en días al año por
    estación para que sean fáciles de interpretar.
\end{itemize}
Cada métrica se reporta junto a su valor sin ataque (las mismas horas limpias) para
distinguir la detección real del azar.

\subsection{Etapa 2: evolución de los detectores (octubre y noviembre de 2026)}
\label{sec:evolucion}

El trabajo avanzó del detector más simple al más complejo, siempre comparando detectores
del mismo tipo, en el mismo escenario y sobre los mismos días («peras con peras», regla
fijada tras la junta del 26 de octubre).

\begin{center}
\small
\renewcommand{\arraystretch}{1.25}
\begin{tabularx}{\linewidth}{@{}>{\raggedright\arraybackslash}p{5.4cm} Y@{}}
  \toprule
  \textbf{Paso} & \textbf{Resultado} \\
  \midrule
  Random Forest, IsolationForest y V4 por hora de red &
    Líneas base. La evaluación por hora de toda la red resultó poco informativa. \\
  LSTM supervisado de red (revisiones 6--8) &
    Tres intentos: ventana de 24~h, entradas relativas al perfil de 14~días y una columna
    por estación que la compara con sus 4~vecinas. Mejora parcial. \\
  Estación-día y regla de vecinas robusta R3 (rev.~11--12) &
    Cambio de unidad de evaluación. R3 es una línea base sin entrenamiento (mediana de
    14~días y dominancia en el vecindario), elegida con 2022. \\
  LSTM por estación (rev.~13) &
    No pasó la elección en 2022 (F1 promedio 0.203 contra 0.330 de R3). Resultado
    negativo. \\
  LSTM local compartido L0 (rev.~14--17) &
    Un solo modelo para las 31~estaciones con la estación y sus 4~vecinas (17~columnas).
    Pasó la elección en 2022 (0.391 contra 0.330) y el examen final de uso único en
    2025--2026 (Tabla~\ref{tab:examen17}). \\
  V4 supervisado y escenarios de una sola estación (rev.~18--19) &
    Comparación justa con el V4: L0 le gana al V4 supervisado compartido en 34 de
    36~combinaciones. Quitar los ataques simultáneos cambia el recall en menos de 0.03. \\
  PARAAD-SA en la RAMA y en Barcelona (rev.~20) &
    L0 obtiene mayor F1 que PARAAD-SA en las 36~combinaciones evaluadas en la RAMA
    (sección~\ref{sec:sota}). \\
  Error fijo en ppb (rev.~21) &
    Segundo tipo de anomalía, pedido por el tutor: «un sensor se equivoca por una
    cantidad, no por un porcentaje». \\
  Línea base anclada en vecinas, semillas y ensamble (LB1, LB2) &
    La ventaja de la base vecinal es estable entre 5~semillas (+0.036 a +0.046), pero no
    se sostuvo en el examen de 2017--2019: resultado negativo de generalización. \\
  Más años de fondo y ensamble Ens-Vb-R (rev.~22) &
    En el examen de despliegue (ago--sep 2026, uso único) superó al ensamble con
    referencia propia: F1 de 0.80 con $\delta$ = 100 y 50\,\%, con 0.2\,\% de falsas
    alarmas. \\
  \textbf{NERAD} = Ens-Vb-R + filtro de vecinas sospechosas (rev.~24) &
    \textbf{Detector final.} No pierde con una sola estación atacada y resiste mejor
    cuando se atacan vecinas. \\
  \bottomrule
\end{tabularx}
\end{center}

\begin{table}[H]
  \centering
  \small
  \caption{Examen final de uso único (revisión~17), F1 con umbral p99, LSTM local L0
    contra la regla de vecinas R3. Modelos congelados, cargados de disco; dos corridas
    idénticas. Con p99.9 y 2025, L0 obtiene 0.756, 0.573 y 0.087 contra 0.632, 0.364 y
    0.019 de R3.}
  \label{tab:examen17}
  \begin{tabular}{@{}l cc cc@{}}
    \toprule
    & \multicolumn{2}{c}{\textbf{2025}} & \multicolumn{2}{c}{\textbf{2026 (ene--jul)}} \\
    \cmidrule(lr){2-3}\cmidrule(lr){4-5}
    \textbf{Anomalía} & L0 & R3 & L0 & R3 \\
    \midrule
    $\pm$100\,\% & 0.457 & 0.475 & 0.540 & 0.476 \\
    $\pm$50\,\%  & 0.445 & 0.391 & 0.502 & 0.390 \\
    $\pm$20\,\%  & 0.249 & 0.120 & 0.242 & 0.092 \\
    \bottomrule
  \end{tabular}
\end{table}

\subsection{El detector final: NERAD}

La Figura~\ref{fig:arquitectura} muestra la arquitectura de NERAD. Sus componentes son:

\begin{enumerate}
  \item \textbf{Entradas.} Para la estación objetivo $s$ y sus 4~vecinas más cercanas, el
    ozono de cada hora y el perfil propio de los últimos 14~días.
  \item \textbf{Filtro vecinal.} Una vecina cuyo desvío respecto a la base vecinal supera
    $\tau = 0.26\,\sigma$ se considera sospechosa y se excluye (se pone en~0 en la entrada).
    El valor de $\tau$ se eligió en 2022.
  \item \textbf{Ventana.} 24~horas $\times$ 17~columnas: el desvío de $s$ contra su base,
    la hora del día (seno y coseno), una máscara de dato presente y el desvío contra las
    vecinas, y lo mismo para cada vecina.
  \item \textbf{Ensamble de 5~LSTM} entrenadas con semillas distintas (42, 1, 2, 3 y 4); su
    promedio da un puntaje por hora $p_s(t)$.
  \item \textbf{Decisión diaria.} Al cierre del día se promedia el puntaje (si hay al menos
    12~horas con dato) y se compara con el umbral p99 calibrado en 2023 sin ataques
    ($S > 0.1508$): si lo supera, la estación-día da alarma.
\end{enumerate}

\begin{figure}[H]
  \centering
  \figurapresentacion{arquitectura}
  \caption{Arquitectura del detector final NERAD: entradas (estación, perfil propio y 4
    vecinas), filtro de vecinas sospechosas, ventana de 24~h $\times$ 17~columnas,
    ensamble de 5~LSTM y decisión diaria por estación.}
  \label{fig:arquitectura}
\end{figure}

% =====================================================================
\section{Resultados principales}
\label{sec:resultados}

\subsection{Elección del punto de operación}

Cada percentil es un umbral distinto fijado en 2023 sin ataques. De p95 a p99.9 el umbral
es más estricto: el recall baja y la precisión sube (Figura~\ref{fig:operacion}). Se
eligió \textbf{p99}: el recall sigue alto (0.80 con $\delta$ = 50\,\% y 0.83 con 20~ppb
en el banco) y las falsas alarmas bajan de 18.5 a 4.7~días al año por estación respecto
de p95. Una falsa alarma cuesta una revisión extra de la estación; por eso se prioriza el
recall sobre la precisión.

\begin{figure}[H]
  \centering
  % Sin \adjustbox: la leyenda va en un parrafo aparte debajo de la grafica y dentro de
  % una caja horizontal se iria a un lado. A su tamano natural (~12.5 cm) ya cabe.
  % Lamina 8: eleccion del punto de operacion del detector final (Ens-Vb-X).
% 2 x 2: columnas = banco | despliegue; filas = recall (arriba) | precision (abajo).
% Banco, p99 y p99.9: results/linea_base/rev24/metricas_banco.csv y falsas_alarmas_banco.csv
%   (Ens-Vb-X0.2596, E1; dias al año = fa x 365). Alineado con la lamina 9 el 2026-11-23.
% Banco, p95 y p98: corrida del registro 24-p98 con p95 agregado (metricas_e1.csv y
%   falsas_alarmas.csv, 2026-11-24), misma mascara que rev24 (23 005 limpias, 1 777 atacadas).
%   Reemplaza a percentiles_banco.csv (mascara comun, no estaba en el repo).
% Despliegue: revision 26, datos validados de SIMAT, unidad A (las mismas estacion-dias
%   del examen original): despliegue_validado_metricas.csv y _falsas_alarmas.csv.
% Umbrales de 2023 sin ataques para p95, p98, p99 y p99.9.
% Eje x: 1 = p95, 2 = p98, 3 = p99, 4 = p99.9 (de menos a mas estricto).
% Eje x: umbral (percentil de 2023 sin ataques). Debajo, en un recuadro aparte (complemento de la
% grafica, no parte de ella): dias con falsa alarma al año por estacion de cada percentil (2026-11-23).
% Tono = tamaño del ataque (claro = chico). Linea solida con circulo = anomalia proporcional
% (delta %); linea punteada con cuadro = error fijo (ppb). Recall arriba, precision abajo.
% Franja crema = p99, el elegido. (Estilos cambiados el 2026-11-23 a pedido del tesista.)
\begin{tikzpicture}[
    recuadroFA/.style={draw=tenue!60, line width=0.4pt, rounded corners=2pt},
    numFA/.style={font=\sffamily\scriptsize, text=tenue},
    textoFA/.style={font=\sffamily\tiny, text=tenue},
  ]
\begin{groupplot}[
    group style={group size=2 by 2, horizontal sep=0.45cm, vertical sep=0.3cm,
                 xticklabels at=edge bottom, yticklabels at=edge left},
    ejes tesis,
    scale only axis, width=5.4cm, height=1.6cm,
    xmin=0.6, xmax=4.4, xtick={1,2,3,4},
    xticklabel style={align=center},
    ymin=0, ymax=1.05, ytick={0,0.5,1},
    every axis plot/.append style={mark size=1.6pt},
    title style={font=\sffamily\small\bfseries, text=noche, yshift=-3pt},
    xlabel style={font=\sffamily\scriptsize, text=texto, yshift=3pt},
    clip=false,
  ]
  % ---------------- Banco · recall ----------------
  \nextgroupplot[title={Banco}, ylabel={Recall},
                 legend to name=leyendaPO, legend columns=4,
                 legend style={/tikz/every even column/.append style={column sep=0.35cm}}]
  \addplot[forget plot, fill=resalte, draw=none, mark=none] coordinates {(2.65,0) (2.65,1.05) (3.35,1.05) (3.35,0)} \closedcycle;  % franja p99
  \addlegendimage{empty legend}\addlegendentry{\textbf{Proporcional} (sólida):}
  \addplot[mag chica, fam pct, met recall] coordinates {(1,0.45) (2,0.29) (3,0.20) (4,0.03)}; \addlegendentry{$\delta$\,20\,\%}
  \addplot[mag media, fam pct, met recall] coordinates {(1,0.92) (2,0.85) (3,0.80) (4,0.40)}; \addlegendentry{$\delta$\,50\,\%}
  \addplot[mag grande, fam pct, met recall] coordinates {(1,0.98) (2,0.95) (3,0.93) (4,0.64)}; \addlegendentry{$\delta$\,100\,\%}
  \addlegendimage{empty legend}\addlegendentry{\textbf{Fijo} (punteada):}
  \addplot[mag chica, fam ppb, met recall] coordinates {(1,0.49) (2,0.29) (3,0.23) (4,0.05)}; \addlegendentry{+10 ppb}
  \addplot[mag media, fam ppb, met recall] coordinates {(1,0.95) (2,0.88) (3,0.83) (4,0.44)}; \addlegendentry{+20 ppb}
  \addplot[mag grande, fam ppb, met recall] coordinates {(1,0.99) (2,0.97) (3,0.96) (4,0.78)}; \addlegendentry{+40 ppb}
  % ---------------- Despliegue · recall ----------------
  \nextgroupplot[title={Despliegue}]
  \addplot[forget plot, fill=resalte, draw=none, mark=none] coordinates {(2.65,0) (2.65,1.05) (3.35,1.05) (3.35,0)} \closedcycle;  % franja p99
  \addplot[mag chica, fam pct, met recall] coordinates {(1,0.45) (2,0.30) (3,0.25) (4,0.03)};  % $\delta$\,20\,\%
  \addplot[mag media, fam pct, met recall] coordinates {(1,0.99) (2,0.97) (3,0.94) (4,0.53)};  % $\delta$\,50\,\%
  \addplot[mag grande, fam pct, met recall] coordinates {(1,1.00) (2,0.99) (3,0.99) (4,0.73)};  % $\delta$\,100\,\%
  \addplot[mag chica, fam ppb, met recall] coordinates {(1,0.55) (2,0.39) (3,0.33) (4,0.06)};  % 10 ppb
  \addplot[mag media, fam ppb, met recall] coordinates {(1,0.97) (2,0.96) (3,0.90) (4,0.53)};  % 20 ppb
  \addplot[mag grande, fam ppb, met recall] coordinates {(1,1.00) (2,1.00) (3,1.00) (4,0.86)};  % 40 ppb
  % ---------------- Banco · precision ----------------
  \nextgroupplot[ylabel={Precisión}, xlabel={Umbral: percentil de 2023 sin ataques},
                 xticklabels={p95, p98, p99, p99.9}]
  % recuadro aparte (complemento, no parte de la grafica): dias con falsa alarma por percentil
  \draw[recuadroFA] ([yshift=-0.93cm]rel axis cs:0,0) rectangle ([yshift=-1.55cm]rel axis cs:1,0);
  \node[numFA] at ([yshift=-1.09cm]axis cs:1,0) {18.5};
  \node[numFA] at ([yshift=-1.09cm]axis cs:2,0) {8.4};
  \node[numFA] at ([yshift=-1.09cm]axis cs:3,0) {4.7};
  \node[numFA] at ([yshift=-1.09cm]axis cs:4,0) {0.7};
  \node[textoFA] at ([yshift=-1.39cm]rel axis cs:0.5,0) {Días con falsa alarma al año por estación};
  \addplot[forget plot, fill=resalte, draw=none, mark=none] coordinates {(2.65,0) (2.65,1.05) (3.35,1.05) (3.35,0)} \closedcycle;  % franja p99
  \node[font=\sffamily\scriptsize\bfseries, text=acento, anchor=north] at (axis cs:3,1.03) {elegido};
  \addplot[mag chica, fam pct, met prec] coordinates {(1,0.06) (2,0.09) (3,0.11) (4,0.13)};  % $\delta$\,20\,\%
  \addplot[mag media, fam pct, met prec] coordinates {(1,0.12) (2,0.22) (3,0.31) (4,0.64)};  % $\delta$\,50\,\%
  \addplot[mag grande, fam pct, met prec] coordinates {(1,0.12) (2,0.22) (3,0.32) (4,0.73)};  % $\delta$\,100\,\%
  \addplot[mag chica, fam ppb, met prec] coordinates {(1,0.07) (2,0.09) (3,0.12) (4,0.17)};  % 10 ppb
  \addplot[mag media, fam ppb, met prec] coordinates {(1,0.12) (2,0.22) (3,0.32) (4,0.65)};  % 20 ppb
  \addplot[mag grande, fam ppb, met prec] coordinates {(1,0.12) (2,0.23) (3,0.33) (4,0.77)};  % 40 ppb
  % ---------------- Despliegue · precision ----------------
  \nextgroupplot[xlabel={Umbral: percentil de 2023 sin ataques},
                 xticklabels={p95, p98, p99, p99.9}]
  % recuadro aparte (complemento, no parte de la grafica): dias con falsa alarma por percentil
  \draw[recuadroFA] ([yshift=-0.93cm]rel axis cs:0,0) rectangle ([yshift=-1.55cm]rel axis cs:1,0);
  \node[numFA] at ([yshift=-1.09cm]axis cs:1,0) {5.4};
  \node[numFA] at ([yshift=-1.09cm]axis cs:2,0) {1.3};
  \node[numFA] at ([yshift=-1.09cm]axis cs:3,0) {1.1};
  \node[numFA] at ([yshift=-1.09cm]axis cs:4,0) {0.0};
  \node[textoFA] at ([yshift=-1.39cm]rel axis cs:0.5,0) {Días con falsa alarma al año por estación};
  \addplot[forget plot, fill=resalte, draw=none, mark=none] coordinates {(2.65,0) (2.65,1.05) (3.35,1.05) (3.35,0)} \closedcycle;  % franja p99
  \addplot[mag chica, fam pct, met prec] coordinates {(1,0.25) (2,0.44) (3,0.46) (4,1.00)};  % $\delta$\,20\,\%
  \addplot[mag media, fam pct, met prec] coordinates {(1,0.39) (2,0.63) (3,0.67) (4,0.96)};  % $\delta$\,50\,\%
  \addplot[mag grande, fam pct, met prec] coordinates {(1,0.35) (2,0.59) (3,0.65) (4,0.96)};  % $\delta$\,100\,\%
  \addplot[mag chica, fam ppb, met prec] coordinates {(1,0.29) (2,0.50) (3,0.52) (4,1.00)};  % 10 ppb
  \addplot[mag media, fam ppb, met prec] coordinates {(1,0.38) (2,0.63) (3,0.67) (4,0.97)};  % 20 ppb
  \addplot[mag grande, fam ppb, met prec] coordinates {(1,0.36) (2,0.59) (3,0.65) (4,0.98)};  % 40 ppb
\end{groupplot}
\end{tikzpicture}

\smallskip
\ref{leyendaPO}   % leyenda comun (se llena en la segunda compilacion)
  \caption{Recall (arriba) y precisión (abajo) de NERAD según el umbral, en el banco de
    pruebas (izquierda) y en el despliegue de agosto--septiembre de 2026 (derecha). Línea
    sólida: anomalía proporcional; punteada: error fijo en ppb; el tono indica la magnitud.
    Debajo de cada gráfica, los días con falsa alarma al año por estación. La franja marca
    p99, el umbral elegido.}
  \label{fig:operacion}
\end{figure}

\subsection{NERAD contra PARAAD-SA en el banco de pruebas}

Las Figuras \ref{fig:comparacion} y~\ref{fig:comparacionppb} y la
Tabla~\ref{tab:banco} comparan a los dos detectores con una sola estación atacada: 1\,777
estación-días atacados y 23\,005 limpios, los mismos para ambos, cada uno con su umbral
p99 de 2023. En el banco, el umbral de PARAAD-SA deja pasar 2.6\,\% de días normales como
alarma (9.6~días al año por estación) y el de NERAD 1.3\,\% (4.7~días).

\begin{figure}[H]
  \centering
  \figurapresentacion{comparacion_pct}
  \caption{Anomalía proporcional en el banco de pruebas (2024 a julio de 2026, 30
    calendarios), umbral p99. Eje vertical: días al año por estación. Por cada magnitud,
    las dos primeras barras son las alarmas (verdaderas en color, falsas en gris; P =
    precisión) y las dos siguientes, los días atacados (detectados en color, perdidos en
    claro; R = recall).}
  \label{fig:comparacion}
\end{figure}

\begin{figure}[H]
  \centering
  \figurapresentacion{comparacion_ppb}
  \caption{Escalón aditivo fijo en el banco de pruebas, umbral p99. Se lee igual que la
    Figura~\ref{fig:comparacion}.}
  \label{fig:comparacionppb}
\end{figure}

\begin{table}[H]
  \centering
  \small
  \caption{Banco de pruebas 2024 a julio de 2026, una estación atacada, umbral p99 de 2023
    sin ataques. Modelos cargados de disco, sin reentrenar.}
  \label{tab:banco}
  \begin{tabular}{@{}l cc cc cc@{}}
    \toprule
    & \multicolumn{2}{c}{\textbf{Recall}} & \multicolumn{2}{c}{\textbf{Precisión}}
    & \multicolumn{2}{c}{\textbf{F1}} \\
    \cmidrule(lr){2-3}\cmidrule(lr){4-5}\cmidrule(lr){6-7}
    \textbf{Anomalía} & NERAD & PARAAD-SA & NERAD & PARAAD-SA & NERAD & PARAAD-SA \\
    \midrule
    $\pm$100\,\%  & 0.93 & 0.68 & 0.32 & 0.16 & 0.48 & 0.27 \\
    $\pm$50\,\%   & 0.80 & 0.36 & 0.31 & 0.10 & 0.44 & 0.15 \\
    $\pm$20\,\%   & 0.20 & 0.09 & 0.11 & 0.03 & 0.14 & 0.04 \\
    $\pm$40~ppb   & 0.96 & 0.92 & 0.33 & 0.21 & 0.49 & 0.34 \\
    $\pm$20~ppb   & 0.83 & 0.57 & 0.32 & 0.14 & 0.46 & 0.23 \\
    $\pm$10~ppb   & 0.23 & 0.10 & 0.12 & 0.03 & 0.16 & 0.05 \\
    \midrule
    \multicolumn{7}{@{}l}{Días con falsa alarma al año por estación: NERAD 4.7;
      PARAAD-SA 9.6.} \\
    \bottomrule
  \end{tabular}
\end{table}

La precisión del banco es baja para los dos detectores porque los días atacados son muy
pocos (0.77\,\% de los estación-días): aun con 1\,\% de falsas alarmas, éstas son del
mismo orden que los ataques.

\subsection{Despliegue (agosto y septiembre de 2026)}

Con los datos validados del SIMAT y el umbral p99, NERAD detecta 99\,\% de los días con
anomalía de 100\,\%, 94\,\% con 50\,\% y 25\,\% con 20\,\%, con precisión de 0.65, 0.67 y
0.46, respectivamente, y 1.1~días con falsa alarma al año por estación. La precisión es
mayor que en el banco porque hubo menos falsas alarmas. Este resultado es descriptivo: la
ventana ya se había usado para examinar el ensamble sin filtro.

\subsection{Comparación con el estado del arte}
\label{sec:sota}

\paragraph{PARAAD-SA en la RAMA.} Se reimplementó el método publicado, se adaptó a las
estaciones de la RAMA, se entrenó sin etiquetas con 2020--2022 y se calibró igual que los
demás. La Tabla~\ref{tab:paraad} muestra el escenario original con $\delta$ = 50\,\%.
PARAAD-SA detecta bien los cambios grandes, pero no muestra una mejora estable frente al
V4 entre años, y el LSTM L0 lo supera en todos los casos.

\begin{table}[H]
  \centering
  \small
  \caption{F1 en la RAMA, $\delta$ = 50\,\%, umbral p99 (revisión~20). Modelos congelados.}
  \label{tab:paraad}
  \begin{tabular}{@{}lccc@{}}
    \toprule
    \textbf{Año} & \textbf{PARAAD-SA} & \textbf{V4} & \textbf{L0} \\
    \midrule
    2024 & 0.128 & 0.101 & 0.388 \\
    2025 & 0.208 & 0.108 & 0.445 \\
    2026 (ene--jul) & 0.140 & 0.193 & 0.502 \\
    \bottomrule
  \end{tabular}
\end{table}

\paragraph{Validación externa en Barcelona.} Se entrenaron L0, V4 y PARAAD-SA con
13~estaciones de Barcelona (ajuste 2014--2016, calibración 2017, examen 2018 y 2019), con
las mismas anomalías proporcionales. L0 obtiene el mayor F1 con $\delta$ = 50 y 20\,\%
en los dos años (Tabla~\ref{tab:barcelona}); con 100\,\% en 2018, PARAAD-SA queda
ligeramente arriba. Todas las inferencias se repitieron dos veces con resultados
idénticos.

\begin{table}[H]
  \centering
  \small
  \caption{F1 en Barcelona, umbral p99. No es una reproducción del artículo original:
    otras estaciones, otro umbral y otra anomalía.}
  \label{tab:barcelona}
  \begin{tabular}{@{}l ccc ccc@{}}
    \toprule
    & \multicolumn{3}{c}{\textbf{2018}} & \multicolumn{3}{c}{\textbf{2019}} \\
    \cmidrule(lr){2-4}\cmidrule(lr){5-7}
    \textbf{Anomalía} & L0 & V4 & PARAAD-SA & L0 & V4 & PARAAD-SA \\
    \midrule
    $\pm$100\,\% & 0.484 & 0.412 & 0.500 & 0.442 & 0.236 & 0.313 \\
    $\pm$50\,\%  & 0.500 & 0.252 & 0.304 & 0.435 & 0.106 & 0.191 \\
    $\pm$20\,\%  & 0.144 & 0.052 & 0.071 & 0.107 & 0.024 & 0.058 \\
    \bottomrule
  \end{tabular}
\end{table}

\subsection{Pruebas de estrés y alcance del detector}
\label{sec:estres}

\paragraph{Ataque coordinado a las vecinas.} Un atacante que altera también a las vecinas
puede engañar a un detector que se apoya en ellas. Con $\delta$ = 50\,\%, el recall de
NERAD sobre la estación atacada es 0.80 sin vecinas alteradas, 0.56 con una, 0.31 con dos
y 0.05 con todas (Figura~\ref{fig:vecinas}). PARAAD-SA, que no usa vecinas, se mantiene
alrededor de 0.31--0.36. El filtro de vecinas obliga al atacante a comprometer 3 o más
estaciones.

\begin{figure}[H]
  \centering
  \figurapresentacion{vecinas_coordinadas}
  \caption{Recall sobre la estación atacada según el número de vecinas también alteradas,
    $\delta$ = 50\,\%, umbral p99, banco de pruebas.}
  \label{fig:vecinas}
\end{figure}

\paragraph{Sesgo sostenido de 8~semanas.} Con +20~ppb, NERAD detecta 0.76 de los días en
las semanas 1--4 contra 0.55 de PARAAD-SA, pero sólo 0.11 en las semanas 5--8, porque su
línea base absorbe el sesgo poco a poco (Figura~\ref{fig:sesgo}). Con +40~ppb el patrón
es el mismo (0.93 y luego 0.20). El tiempo máximo de detección es de unas 4~semanas.
También se identificó un riesgo operativo: al terminar un sesgo largo hay unas 7~semanas
de alarmas de rebote, que aparecerían igual tras una recalibración legítima del sensor;
el procedimiento sugerido es reiniciar la línea base después de cada mantenimiento.

\begin{figure}[H]
  \centering
  \figurapresentacion{sesgo_sostenido}
  \caption{Recall por semana ante un sesgo sostenido de +20~ppb (izquierda) y +40~ppb
    (derecha), umbral p99; 240~episodios aislados de sus vecinas.}
  \label{fig:sesgo}
\end{figure}

\paragraph{Fusión NERAD + PARAAD-SA.} Como los dos detectores fallan en casos distintos,
se probó combinarlos (revisión~25). Con todas las vecinas comprometidas, la fusión sube el
recall de 0.05 a 0.26; pero con una sola estación atacada, que es la amenaza principal,
su precisión cae de 0.31 a 0.18 y pierde en F1 ($-0.110$), porque el umbral de PARAAD-SA
produce 2.6\,\% de falsas alarmas reales en lugar del 1\,\% nominal. Por eso la fusión no
forma parte del sistema propuesto y queda como trabajo futuro, con PARAAD-SA recalibrado
periódicamente.

\subsection{Exploración con imágenes sintéticas TINTO}
\label{sec:tinto}

TINTO convierte datos tabulares en imágenes: cada estación-día se representa como una
imagen de $32\times32$ píxeles en la que la posición de cada dato se fija con un mapa de
componentes principales. Sobre esas imágenes se entrenan redes convolucionales. Se
cerraron dos pilotos (T1 y T2) y se exploraron dos variantes (Figura~\ref{fig:tinto}):
un autoencoder que mira un solo día y una CNN temporal que mira los 7~días previos.

\begin{figure}[H]
  \centering
  \figurapresentacion{tinto}
  \caption{Exploración con imágenes TINTO, umbral p99, sobre los mismos 175~estación-días
    para los dos modelos. «+100\,\%» se refiere sólo al sentido de aumento.}
  \label{fig:tinto}
\end{figure}

\paragraph{Comparación pareada con el detector final (11 de noviembre).} Con un modelo
TINTO congelado (CNN + GRU diaria, 11\,104~parámetros) y NERAD sin cambios, sobre los
mismos días de 2025 y un calendario de ataque (Tabla~\ref{tab:tinto}):

\begin{table}[H]
  \centering
  \small
  \caption{TINTO contra NERAD, 2025, umbral p99, 42~estación-días atacados por escenario y
    3\,067 limpios. Falsas alarmas en los días originales: 73 de TINTO y 21 de NERAD en
    3\,109~estación-días.}
  \label{tab:tinto}
  \begin{tabular}{@{}l cc cc@{}}
    \toprule
    & \multicolumn{2}{c}{\textbf{Recall}} & \multicolumn{2}{c}{\textbf{Precisión}} \\
    \cmidrule(lr){2-3}\cmidrule(lr){4-5}
    \textbf{Anomalía} & TINTO & NERAD & TINTO & NERAD \\
    \midrule
    $\times 2$   & 0.71 & 1.00 & 0.28 & 0.58 \\
    $\times 1.5$ & 0.40 & 0.90 & 0.19 & 0.59 \\
    +20~ppb      & 0.33 & 0.98 & 0.16 & 0.61 \\
    +40~ppb      & 0.60 & 1.00 & 0.24 & 0.60 \\
    \bottomrule
  \end{tabular}
\end{table}

En este piloto TINTO no supera al detector final: usa muchos menos parámetros y no
necesita vecinas, pero tiene más falsas alarmas, peor detección y menor cobertura (3\,109
contra 8\,822~días elegibles). Es una muestra pequeña con un solo calendario y un año ya
usado, así que no es un veredicto definitivo. La línea sigue en curso.

\subsection{Lectura de los resultados}

\begin{itemize}
  \item NERAD cumple la meta de detección para anomalías de 50\,\% o más y para escalones
    de 20~ppb o más, con menos de la mitad de falsas alarmas que PARAAD-SA.
  \item Las anomalías chicas (20\,\% o 10~ppb) siguen siendo difíciles para todos los
    detectores probados: están dentro de la variación natural del ozono.
  \item Ningún detector cubre por sí solo las tres amenazas (un nodo, vecindario
    comprometido y sesgo lento). NERAD es el mejor en la amenaza principal y sus dos
    puntos ciegos quedan medidos y declarados.
  \item Respecto a la hipótesis: los resultados la apoyan para anomalías moderadas y
    grandes; para anomalías pequeñas la evidencia todavía no es suficiente.
\end{itemize}

% =====================================================================
\section{Rigor metodológico}

\begin{itemize}
  \item \textbf{Registro previo:} cada experimento tiene definiciones, métricas, regla de
    decisión y predicciones escritas y fechadas antes de ver resultados. Las decisiones
    tomadas después de ver resultados se declaran como tales.
  \item \textbf{Separación temporal:} nada se ajusta mirando los datos de prueba; los
    exámenes finales se usan una sola vez (2025--2026 para L0, agosto--septiembre de 2026
    para el ensamble).
  \item \textbf{Reproducibilidad:} semilla fija (42) antes de cada entrenamiento, cada
    resultado se corre dos veces y se comparan las salidas, y los modelos y puntajes
    llevan huellas SHA-256. Se indica siempre si un número viene de un modelo recién
    entrenado o cargado de disco.
  \item \textbf{Comparaciones justas:} detectores del mismo tipo, en el mismo escenario,
    con los mismos calendarios y días.
  \item \textbf{Intervalos de confianza:} las diferencias de F1 entre detectores se
    reportan con intervalos por \emph{bootstrap} de bloques estación~$\times$~mes.
  \item \textbf{Resultados negativos documentados:} la codificación alternativa (etapa~1),
    el LSTM por estación, la acumulación entre días, la mediana en lugar de la media y la
    generalización de la base vecinal a 2017--2019.
  \item \textbf{Cómputo:} los entrenamientos largos se ejecutaron en Google Colab Pro sólo
    con procesador (para mantener el determinismo) y en una PC local.
\end{itemize}

% =====================================================================
\section{Limitaciones reconocidas}

\begin{itemize}
  \item Las anomalías son sintéticas, inyectadas sobre datos reales; no se dispone de
    anomalías reales etiquetadas.
  \item El despliegue cubre sólo dos meses. Además, el filtro de vecinas se diseñó después
    del examen de despliegue del ensamble, así que esa ventana ya no es un examen limpio
    para NERAD y su resultado se presenta como descriptivo.
  \item La ventaja de la base anclada en vecinas depende del periodo: no se sostuvo en el
    examen de 2017--2019.
  \item Puntos ciegos declarados: vecindario completo comprometido y sesgo sostenido de más
    de unas 4~semanas.
  \item Las estaciones de la RAMA están relativamente lejos entre sí, lo que limita el
    contraste con las vecinas.
  \item La comparación con PARAAD-SA usa una reimplementación adaptada a la RAMA, no el
    código original de los autores.
\end{itemize}

% =====================================================================
\section{Productos generados}

\begin{itemize}
  \item Repositorio de la tesis con el código (ingesta, generación de anomalías,
    detectores y evaluación), pruebas automáticas y resultados versionados.
  \item Protocolo experimental de la etapa~2 con sus revisiones fechadas, registros
    previos de cada experimento, bitácora de decisiones y registro de progreso por sesión.
  \item Cuadernos (\emph{notebooks}) que explican desde cero el entrenamiento y la
    evaluación del LSTM.
  \item Base de referencias verificadas y análisis del estado del arte (PARAAD-SA,
    STC-RAAD y comparadores).
  \item Presentación de avances en Beamer con gráficas generadas en pgfplots a partir de
    los resultados, y un informe de la comparación L0, V4 y PARAAD-SA.
\end{itemize}

\section{Otras actividades académicas}

\porllenar{Cursos, seminarios, estancias, congresos o coloquios del periodo, por llenar.}

% =====================================================================
\section{Estado de la tesis y plan de trabajo}

\begin{center}
\small
\begin{tabularx}{\linewidth}{@{}l l Y@{}}
  \toprule
  \textbf{Capítulo} & \textbf{Estado} & \textbf{Comentario} \\
  \midrule
  1. Introducción & Borrador & Problema, hipótesis y objetivos; actualizar con la etapa~2. \\
  2. Marco teórico & Por redactar & LoRaWAN y AES-CTR, LSTM, detección de anomalías. \\
  3. Estado del arte & Material reunido & Análisis de PARAAD-SA, STC-RAAD y FedWLA listo. \\
  4. Metodología & Material reunido & Protocolo y bitácora completos. \\
  5. Resultados & Experimentos cerrados & Falta redactar; gráficas ya hechas. \\
  6. Conclusiones & Por redactar & --- \\
  \bottomrule
\end{tabularx}
\end{center}

\begin{figure}[H]
  \centering
  \figurapresentacion{trabajo_futuro}
  \caption{Plan de trabajo: lo hecho y lo que falta.}
  \label{fig:futuro}
\end{figure}

\paragraph{Plan propuesto} \porllenar{fechas por confirmar con el tutor}:
\begin{enumerate}
  \item Cerrar la exploración con TINTO con una comparación más amplia contra NERAD sobre
    los mismos días.
  \item Redactar los capítulos 2 a 6 de la tesis.
  \item Preparar y enviar un artículo con los resultados de NERAD.
\end{enumerate}

\paragraph{Trabajo futuro identificado} (fuera del alcance actual):
\begin{itemize}
  \item Fusión con un detector sin vecinas (PARAAD-SA recalibrado periódicamente) para
    cubrir el ataque coordinado y el sesgo de largo plazo.
  \item Comparación directa con STC-RAAD, que también usa a las vecinas.
  \item Un atacante que optimiza su ataque contra el detector (\emph{red team}).
  \item Extensión a otros contaminantes (PM$_{\text{10}}$, PM$_{\text{2.5}}$).
\end{itemize}

% =====================================================================
\section{Conclusiones del periodo}

En este periodo la tesis pasó de caracterizar un ataque específico a proponer un detector
general de estaciones anómalas. La etapa~1 dejó una contribución cerrada: qué posiciones
de bit de LoRaWAN~1.1 producen daño normativo y cuáles detecta un LSTM. La etapa~2 produjo
NERAD, que en un banco de 2024 a julio de 2026 detecta 80\,\% de las anomalías de 50\,\% y
83\,\% de los escalones de 20~ppb con 4.7~días de falsa alarma al año por estación, y que
supera a un método publicado reciente con la mitad de falsas alarmas. Los límites del
detector están medidos y declarados. Con los experimentos principales cerrados, el
siguiente semestre se dedica sobre todo a la redacción de la tesis y de un artículo.

\printbibliography[heading=bibintoc, title={Referencias}]

\end{document}
